\chapter{From Neural Architecture to Robot Control}
\label{ch:neural-to-robot}

\epigraph{%
  If we truly understood how humans move,
  we could build robots that move like humans.
  We cannot. This tells us something.
}{}

\section{Lessons from Biology for Robot Design}

The four volumes of theory converge on practical design principles for robotic
control systems:

\begin{insight}{Bio-Inspired Control Design Principles}{bio-robot-design}
  \begin{enumerate}
    \item \textbf{Exploit passive dynamics} (Vol~II, Vol~IV Ch~7): design the mechanical
          system so that natural dynamics do useful work. Use compliance, not rigidity.
    \item \textbf{Use shallow-wide controllers} (Vol~IV Ch~4): for real-time control,
          prefer 2--3 layer networks with thousands of neurons over deep networks with
          hundreds of layers.
    \item \textbf{Employ synergy-based dimensionality reduction} (Vol~IV Ch~4): define
          4--8 actuation synergies and control in synergy space, not muscle space.
    \item \textbf{Implement forward models} (Vol~IV Ch~6): use the Smith predictor
          architecture to compensate for sensor/actuator delays.
    \item \textbf{Design for contraction} (Vol~I Ch~4): ensure the closed-loop
          system is contracting, guaranteeing stability and robustness.
    \item \textbf{Optimize trajectories, not setpoints} (Vol~II): design
          trajectory-centric controllers with funnel guarantees.
    \item \textbf{Accept variability} (Vol~IV Ch~1): allow task-irrelevant
          variability. Don't over-constrain the controller.
  \end{enumerate}
\end{insight}

\section{Cerebellar-Inspired Adaptive Control}

The cerebellum's architecture suggests a specific adaptive controller:
\begin{itemize}
  \item \textbf{Granule cells} $\to$ basis function expansion (like kernel methods)
  \item \textbf{Purkinje cells} $\to$ linear combination of basis functions
  \item \textbf{Climbing fiber error} $\to$ supervised learning signal
\end{itemize}

This maps directly to:
\begin{equation}
  \control_{\text{adaptive}}(\state) = \hat{W}^T \bm{\phi}(\state),
  \quad \dot{\hat{W}} = -\Gamma \bm{\phi}(\state) \bm{e}^T,
  \label{eq:ch11:cerebellar-control}
\end{equation}
where $\bm{\phi}(\state)$ are radial basis functions (granule cell expansion),
$\hat{W}$ are adaptive weights (parallel fiber--Purkinje cell synapses), and
$\bm{e}$ is the tracking error (climbing fiber signal).

\section{CPG-Based Locomotion Controllers}

Following Chapter~8, CPG-based locomotion controllers use coupled oscillators to
generate walking/running patterns. The CPG provides the \emph{timing} and the
musculoskeletal model provides the \emph{mechanics}:

\begin{lstlisting}[language=Python, caption={Bio-inspired locomotion controller architecture.}]
import numpy as np

class BioInspiredLocomotionController:
    """Hierarchical control architecture inspired by
    biological motor control."""

    def __init__(
        self,
        n_joints: int,
        n_synergies: int = 4,
        cpg_freq: float = 1.0
    ):
        self.n_joints = n_joints
        self.n_synergies = n_synergies
        self.cpg_freq = cpg_freq

        # Synergy matrix (learned from motion data)
        self.W = np.random.randn(n_joints, n_synergies) * 0.1

        # CPG phases for each synergy
        self.phases = np.linspace(
            0, 2 * np.pi * (1 - 1/n_synergies), n_synergies
        )

        # Feedback gains (shallow: single layer)
        self.K_fb = np.zeros((n_synergies, 2 * n_joints))

    def cpg_output(self, t: float) -> np.ndarray:
        """Generate rhythmic synergy activations from CPG."""
        activations = np.array([
            max(0.0, np.sin(2 * np.pi * self.cpg_freq * t + phi))
            for phi in self.phases
        ])
        return activations

    def feedforward(self, t: float) -> np.ndarray:
        """Feedforward joint torques from CPG + synergies."""
        c = self.cpg_output(t)
        return self.W @ c

    def feedback(self, state_error: np.ndarray) -> np.ndarray:
        """Feedback correction (shallow: one layer)."""
        c_fb = self.K_fb @ state_error
        return self.W @ c_fb

    def control(self, t: float,
                state_error: np.ndarray) -> np.ndarray:
        """Total control: feedforward (CPG) + feedback."""
        return self.feedforward(t) + self.feedback(state_error)


# Example: 6-joint biped with 4 synergies
controller = BioInspiredLocomotionController(
    n_joints=6, n_synergies=4, cpg_freq=1.2
)

# Generate one cycle of control signals
t = np.linspace(0, 1.0/1.2, 200)
torques = np.array([controller.feedforward(ti) for ti in t])

print(f"Bio-inspired controller: {controller.n_joints} joints, "
      f"{controller.n_synergies} synergies")
print(f"Control dimensionality reduced from {controller.n_joints} "
      f"to {controller.n_synergies}")
print(f"Peak torque magnitude: {np.max(np.abs(torques)):.3f}")
\end{lstlisting}

\section{Whole-Body Control Architectures}

For humanoid robots, the bio-inspired architecture suggests:

\begin{enumerate}
  \item \textbf{Level 0 (Spinal)}: joint-level impedance control with reflex-like
        feedback ($\sim$1~ms loop)
  \item \textbf{Level 1 (CPG)}: rhythmic pattern generation for locomotion
        ($\sim$10~ms loop)
  \item \textbf{Level 2 (Cerebellar)}: adaptive feedforward based on forward model
        predictions ($\sim$50~ms loop)
  \item \textbf{Level 3 (Cortical)}: trajectory planning and task-level control
        ($\sim$100~ms loop)
\end{enumerate}

Each level is \textbf{shallow} (1--2 layers) but \textbf{wide} (many parallel channels),
operating at progressively longer timescales.

\section{The Connection to UpstreamDrift}

Volume~V provides the practical tools to implement and test these controllers using
the UpstreamDrift simulation platform. The connection points:
\begin{itemize}
  \item \textbf{MuJoCo}: test CPG-based locomotion controllers on musculoskeletal models
  \item \textbf{Drake}: implement trajectory optimization for trajectory-centric control
  \item \textbf{Pinocchio}: compute forward/inverse dynamics efficiently in real-time
\end{itemize}

\section*{Exercises}

\begin{enumerate}
  \item \textbf{Controller Architecture.}
        Using the bio-inspired controller code, vary the number of synergies and
        measure the resulting control signal dimensionality and smoothness.

  \item \textbf{Cerebellar Adaptive Control.}
        Implement the cerebellar-inspired adaptive controller for a 2-DOF arm and
        show that it learns an inverse model over 50 trials.

  \item \textbf{Hierarchical Design.}
        Design a 3-level hierarchical controller for a simulated biped, with
        impedance control at Level 0, CPG at Level 1, and trajectory correction at Level 2.

  \item \textbf{(Programming.)} Implement the full bio-inspired controller for a
        MuJoCo walking model and benchmark against a standard PD controller.
\end{enumerate}
